\chapter*{Glosario}
\addcontentsline{toc}{chapter}{Glosario}

\begin{description}

\item[Análisis léxico:]
Fase del procesamiento de un programa que transforma el texto fuente en una secuencia de unidades léxicas llamadas tokens.

\item[Análisis sintáctico (Parsing):]
Proceso mediante el cual se verifica si una secuencia de tokens cumple las reglas de una gramática y se construye una representación estructural del programa.

\item[AST (Árbol de Sintaxis Abstracta):]
Estructura jerárquica que representa la organización sintáctica y semántica esencial de un programa, omitiendo detalles superficiales de la sintaxis concreta. Sirve como base para el análisis semántico, la interpretación y la generación de código.

\item[CFG (Gramática Libre de Contexto):]
Formalismo matemático compuesto por terminales, no terminales y reglas de producción que define la sintaxis de un lenguaje de programación y permite describir su estructura de manera jerárquica.

\item[Gramática:]
Conjunto de reglas formales que especifica cómo se construyen las cadenas válidas de un lenguaje.

\item[Intérprete:]
Sistema que ejecuta un programa analizando y evaluando sus instrucciones sin generar previamente código máquina.

\item[Lexer (Analizador léxico):]
Componente que reconoce tokens a partir del flujo de caracteres del código fuente.

\item[Parser (Analizador sintáctico):]
Componente que implementa el análisis sintáctico y genera estructuras como árboles de sintaxis.

\item[Semántica:]
Conjunto de reglas que definen el significado y comportamiento de las construcciones de un lenguaje más allá de su forma sintáctica.

\item[Token:]
Unidad léxica mínima con significado sintáctico, como identificadores, palabras reservadas, operadores o símbolos.

\end{description}
